\documentclass[12pt]{article}
\usepackage[ukrainian,shorthands=off]{babel}   % shorthands=off: символ " не активний (потрібно для міток "..." у TikZ всередині \zadtask)
\usepackage{fontspec}
% --- НАЛАШТУВАННЯ ШРИФТІВ ---
% Шрифти Schoolbook-*.otf лежать у корені проєкту. Шлях підбирається автоматично:
% ../ (компіляція з папки теми), ./ (корінь проєкту), ../../ (підпапки теми 28); інакше -- TeX Gyre Schola.
\newcommand{\nmtSchoolbook}[1]{\setmainfont{Schoolbook-Regular.otf}[
    Path = #1,
    BoldFont = Schoolbook-Bold.otf,
    ItalicFont = Schoolbook-Italic.otf,
    BoldItalicFont = Schoolbook-BoldItalic.otf
]}
\IfFileExists{../Schoolbook-Regular.otf}{\nmtSchoolbook{../}}{%
 \IfFileExists{./Schoolbook-Regular.otf}{\nmtSchoolbook{./}}{%
  \IfFileExists{../../Schoolbook-Regular.otf}{\nmtSchoolbook{../../}}{%
   \setmainfont{texgyreschola}[Extension=.otf, UprightFont=*-regular, BoldFont=*-bold, ItalicFont=*-italic, BoldItalicFont=*-bolditalic]}}}
\usepackage[italic]{mathastext}
\usepackage[a4paper,margin=1.8cm,bottom=2cm,top=2cm]{geometry}
\usepackage{amsmath,amssymb}
\usepackage{tikz}
\usepackage{pgfplots}
\pgfplotsset{compat=1.18}
\usepgfplotslibrary{fillbetween}
\usetikzlibrary{calc,patterns,angles,quotes,intersections,babel,3d,shapes.symbols,shapes.geometric,decorations.pathreplacing,shadings,arrows.meta,hobby}
\usepackage{xcolor,array,fancyhdr,enumitem,multicol}
\usepackage{colortbl,multirow,diagbox,needspace}
\usepackage{tcolorbox}
\tcbuselibrary{skins,breakable}
\usepackage[hidelinks]{hyperref}
\usepackage[firstpage=false, text=@pvtr2525, scale=0.6, color=gray!12]{draftwatermark}

% --- АВТОСТИСКАННЯ: рисунок/сітка, ширша за колонку, масштабується до ширини колонки ---
\newsavebox{\nmtfitbox}
\newenvironment{nmtfit}{\begin{lrbox}{\nmtfitbox}}{\end{lrbox}%
  \ifdim\wd\nmtfitbox>\linewidth \resizebox{\linewidth}{!}{\usebox{\nmtfitbox}}\else\usebox{\nmtfitbox}\fi}
\let\nmtIncludegraphicsOrig\includegraphics
\renewcommand{\includegraphics}[2][]{\begin{nmtfit}\nmtIncludegraphicsOrig[#1]{#2}\end{nmtfit}}


\definecolor{mainGreen}{RGB}{34, 120, 64}
\definecolor{yearOrange}{RGB}{220, 100, 30}
\definecolor{yearcolor}{RGB}{220, 100, 30}
\definecolor{titleBg}{RGB}{225, 240, 220}
\definecolor{instrBg}{RGB}{255, 240, 220}
\definecolor{instrBorder}{RGB}{220, 150, 80}
\definecolor{headerblue}{RGB}{0, 102, 204}   % використовується в рисунках старої бази

\setlength{\parindent}{0pt}
\emergencystretch=1.5em

\pagestyle{fancy}
\fancyhf{}
\renewcommand{\headrulewidth}{0.4pt}
\renewcommand{\footrulewidth}{0.4pt}
\fancyhead[C]{\small\color{gray!80} База завдань НМТ 2022--2026 \textendash{} Тема 28b}
\fancyhead[R]{\small\color{mainGreen}\textbf{\href{https://t.me/pvtr2525}{@pvtr2525}}}
\fancyfoot[L]{\small\color{gray!80}\href{https://t.me/pvtr2525}{ТГ: @pvtr2525}}
\fancyfoot[C]{\small\color{gray!80}\thepage}
\fancyfoot[R]{\small\color{gray!80}НМТ 2022--2026}

% --- ТАБЛИЦІ ВІДПОВІДЕЙ ---
% ширина клітинки = (ширина рядка - відступи) / 5: на всю сторінку це ~3 см, у minipage поруч із рисунком таблиця стискається
\newlength{\nmtcellw}
\newcommand{\nmtSetCell}{\setlength{\nmtcellw}{\dimexpr(\linewidth-12\tabcolsep-6\arrayrulewidth)/5\relax}}
\newcommand{\answerTable}[5]{%
\par\nopagebreak\vspace{0.15cm}
\noindent\nmtSetCell
\begin{tabular}{|*{5}{>{\centering\arraybackslash}m{\nmtcellw}|}}
\hline
\rule[-0.15cm]{0pt}{0.6cm}\textbf{А} & \rule[-0.15cm]{0pt}{0.6cm}\textbf{Б} & \rule[-0.15cm]{0pt}{0.6cm}\textbf{В} & \rule[-0.15cm]{0pt}{0.6cm}\textbf{Г} & \rule[-0.15cm]{0pt}{0.6cm}\textbf{Д} \\
\hline
\rule[-0.2cm]{0pt}{0.75cm}#1 & \rule[-0.2cm]{0pt}{0.75cm}#2 & \rule[-0.2cm]{0pt}{0.75cm}#3 & \rule[-0.2cm]{0pt}{0.75cm}#4 & \rule[-0.2cm]{0pt}{0.75cm}#5 \\
\hline
\end{tabular}
\par\vspace{0.25cm}
}
\newcommand{\answerTableTall}[5]{%
\par\nopagebreak\vspace{0.15cm}
\noindent\nmtSetCell
\begin{tabular}{|*{5}{>{\centering\arraybackslash}m{\nmtcellw}|}}
\hline
\rule[-0.2cm]{0pt}{0.7cm}\textbf{А} & \rule[-0.2cm]{0pt}{0.7cm}\textbf{Б} & \rule[-0.2cm]{0pt}{0.7cm}\textbf{В} & \rule[-0.2cm]{0pt}{0.7cm}\textbf{Г} & \rule[-0.2cm]{0pt}{0.7cm}\textbf{Д} \\
\hline
\rule[-0.45cm]{0pt}{1.2cm}#1 & \rule[-0.45cm]{0pt}{1.2cm}#2 & \rule[-0.45cm]{0pt}{1.2cm}#3 & \rule[-0.45cm]{0pt}{1.2cm}#4 & \rule[-0.45cm]{0pt}{1.2cm}#5 \\
\hline
\end{tabular}
\par\vspace{0.25cm}
}
\newcommand{\instructionBox}[1]{%
\par\vspace{0.3cm}
\begin{tcolorbox}[colback=instrBg, colframe=instrBorder, boxrule=0.6pt, arc=2pt,
    left=8pt, right=8pt, top=4pt, bottom=4pt]
\centering\small\bfseries #1
\end{tcolorbox}
\par\vspace{0.15cm}
}
\newcommand{\sectionTitle}[1]{%
\par\vspace{0.4cm}
\begin{tcolorbox}[colback=titleBg, colframe=titleBg, boxrule=0pt, arc=8pt,
    halign=center, fontupper=\Large\bfseries\color{mainGreen}]
\hyphenpenalty=10000 \exhyphenpenalty=10000 #1
\end{tcolorbox}
\par\vspace{0.3cm}
}
\newcommand{\task}[2]{\noindent\textbf{#1.}\ \ #2\par\vspace{0.2cm}}
% --- НМТ-СТИЛЬ ПОЛЕ ВІДПОВІДІ ---
\newcommand{\nmtAnswerBox}{%
\par\nopagebreak\vspace{0.2cm}\noindent
Відповідь:\ \
\begin{tabular}{@{}*{4}{|>{\centering\arraybackslash}p{0.8cm}}|@{\hspace{0.1cm}\textbf{,}\hspace{0.1cm}}*{3}{|>{\centering\arraybackslash}p{0.8cm}}|@{}}
\hline
\rule[-0.2cm]{0pt}{0.7cm} & & & & & & \\
\hline
\end{tabular}
\par\vspace{0.3cm}
}
% --- СІТКА ДЛЯ МАТЧИНГУ ---
\newsavebox{\nmtgridbox}
\newcommand{\matchingGrid}{%
    \sbox{\nmtgridbox}{\matchingGridRaw}%
    \ifdim\wd\nmtgridbox>\linewidth \resizebox{\linewidth}{!}{\usebox{\nmtgridbox}}\else\usebox{\nmtgridbox}\fi}
\newcommand{\matchingGridRaw}{%
    \begingroup
    \renewcommand{\arraystretch}{1.15}
    \setlength{\tabcolsep}{4pt}
    \footnotesize
    \begin{tabular}{r|c|c|c|c|c|}
         \multicolumn{1}{c}{} & \multicolumn{1}{c}{\textbf{А}} & \multicolumn{1}{c}{\textbf{Б}} & \multicolumn{1}{c}{\textbf{В}} & \multicolumn{1}{c}{\textbf{Г}} & \multicolumn{1}{c}{\textbf{Д}} \\ \cline{2-6}
         \textbf{1} & \phantom{X} & \phantom{X} & \phantom{X} & \phantom{X} & \phantom{X} \\ \cline{2-6}
         \textbf{2} & & & & & \\ \cline{2-6}
         \textbf{3} & & & & & \\ \cline{2-6}
    \end{tabular}
    \endgroup
}
% --- ДОДАТКОВІ МАКРОСИ ЗБІРНИКА ---
\newcommand{\taskBlock}[1]{%
\par\vspace{0.45cm}
\noindent{\color{mainGreen}\rule{\linewidth}{0.8pt}}\par\vspace{0.12cm}
\noindent{\small\bfseries\color{mainGreen}#1}\par\vspace{0.25cm}
}
\newcounter{zad}
\newcommand{\zadtask}[1]{\stepcounter{zad}\noindent\textbf{\thezad.}\ \ #1\par\nopagebreak\vspace{0.2cm}}
\newcommand{\zadnum}{\stepcounter{zad}\textbf{\thezad.}\ \ }
\newcounter{ans}
\newcommand{\ansitem}[1]{\stepcounter{ans}\noindent\textbf{\theans.}\ #1\par\vspace{0.07cm}}
% мітка року: нерозривна, притиснута праворуч; якщо не вміщається -- праворуч на новому рядку
\newcommand{\nmtyear}[1]{\unskip\nobreak\finalhyphendemerits=0 \hfill\penalty100\hskip1em\hbox{}\nobreak\hfill{\small\color{yearOrange}\mbox{\rule{0pt}{2.3ex}(НМТ~#1)}}}
% --- МАКРОС КОРОТКОГО РОЗВ'ЯЗКУ ---
\newcommand{\solution}[1]{%
\par\vspace{0.12cm}
\begin{tcolorbox}[colback=titleBg!45, colframe=mainGreen!55, boxrule=0.4pt, arc=2pt,
    left=6pt, right=6pt, top=3pt, bottom=3pt, breakable]
\footnotesize\textbf{Короткий розв'язок.}\ #1
\end{tcolorbox}
\par\vspace{0.12cm}
}
% Розв'язки й відповіді (є до завдань НМТ-2026): \showsolutionstrue -- показати, \showsolutionsfalse -- сховати
\newif\ifshowsolutions
\showsolutionsfalse
% --- ЗАГОЛОВКИ З ПУНКТАМИ КЛІКАБЕЛЬНОГО ЗМІСТУ ---
\setcounter{tocdepth}{2}
\newcommand{\chapterTitle}[1]{%
\clearpage\phantomsection\addcontentsline{toc}{section}{#1}%
\sectionTitle{#1}}
\newcommand{\typeTitle}[1]{%
\needspace{6\baselineskip}%
\phantomsection\addcontentsline{toc}{subsection}{\quad #1}%
\par\vspace{0.3cm}
\noindent{\large\bfseries\color{yearOrange}#1}\par\vspace{0.08cm}
\noindent{\color{yearOrange}\rule{\linewidth}{0.4pt}}\par\nopagebreak\vspace{0.2cm}}
\raggedbottom
% усередині одного завдання сторінка не розривається: нерозривний вертикальний проміжок
% і нерозривна межа абзаців (умова ніколи не відділяється від варіантів, рисунка чи сітки)
\newcommand{\nbvspace}[1]{\nopagebreak\vspace{#1}}
\newcommand{\nmtnobreak}{\ifvmode\penalty10000 \fi}
% --- ЄДИНИЙ МАКЕТ ЗАВДАНЬ НА ВІДПОВІДНІСТЬ ---
% мітка в колонці сталої ширини, текст -- у parbox: продовження довгого рядка
% вирівнюється під текстом, а не під міткою; відступ між пунктами однаковий скрізь
\newlength{\nmtMatchLab}\setlength{\nmtMatchLab}{1.6em}
\newlength{\nmtMatchSep}\setlength{\nmtMatchSep}{0.28cm}
\newcommand{\matchHead}[1]{\par\noindent\textit{#1}\par\nopagebreak\vspace{0.2cm}}
\newcommand{\matchItem}[2]{\par\noindent\makebox[\nmtMatchLab][l]{\textbf{#1}}%
\parbox[t]{\dimexpr\linewidth-\nmtMatchLab\relax}{\raggedright #2}\par\nopagebreak\vspace{\nmtMatchSep}}
\newcommand{\ansTheme}[1]{\par\vspace{0.25cm}\noindent{\bfseries\color{mainGreen}#1}\par\vspace{0.12cm}}
\newcommand{\ansType}[1]{\par\vspace{0.1cm}\noindent{\itshape\color{yearOrange}#1}\par\vspace{0.08cm}}

\begin{document}
\begingroup\setcounter{zad}{0}
% --- сумісність зі старою базою (діє, якщо тема не має власного визначення) ---
\providecommand{\trigStyle}[1]{#1}
\providecommand{\answerTableSmall}[5]{%
\begingroup\setlength{\tabcolsep}{2pt}\small\nmtSetCell
\begin{tabular}{|*{5}{>{\centering\arraybackslash}m{\nmtcellw}|}}
\hline
\rule[-0.2cm]{0pt}{0.6cm}\textbf{А} & \textbf{Б} & \textbf{В} & \textbf{Г} & \textbf{Д} \\
\hline
\rule[-0.4cm]{0pt}{0.9cm}#1 & \rule[-0.4cm]{0pt}{0.9cm}#2 & \rule[-0.4cm]{0pt}{0.9cm}#3 & \rule[-0.4cm]{0pt}{0.9cm}#4 & \rule[-0.4cm]{0pt}{0.9cm}#5 \\
\hline
\end{tabular}\endgroup}
\providecommand{\matchingLayout}[3]{%
    \noindent
    \begin{minipage}[t]{0.36\textwidth}\vspace{0pt}\raggedright
        #1
    \end{minipage}%
    \hfill
    \begin{minipage}[t]{0.43\textwidth}\vspace{0pt}\raggedright
        #2
    \end{minipage}%
    \hfill
    \begin{minipage}[t]{0.19\textwidth}
        \vspace{0pt}
        \begin{flushright}
        #3
        \end{flushright}
    \end{minipage}}
\providecommand{\matchingLayoutBottom}[2]{%
    \noindent
    \begin{minipage}[t]{0.48\textwidth}\vspace{0pt}\raggedright
        #1
    \end{minipage}%
    \hfill
    \begin{minipage}[t]{0.48\textwidth}\vspace{0pt}\raggedright
        #2
    \end{minipage}
    \vspace{0.5cm}
    \noindent
    \matchingGrid}
\providecommand{\answerListVertical}[5]{%
    \par\nopagebreak\vspace{0.2cm}
    \begin{itemize}[itemsep=0.4cm, leftmargin=1.5cm, labelsep=0.5cm]
        \item[\textbf{А}] #1
        \item[\textbf{Б}] #2
        \item[\textbf{В}] #3
        \item[\textbf{Г}] #4
        \item[\textbf{Д}] #5
    \end{itemize}
    \vspace{0.2cm}}

\chapterTitle{Тема 28b. Показникова функція і показникові вирази}
\noindent{\small\color{gray!80} Розділ 3. Функції \quad\textbullet\quad Усього завдань: 31 (НМТ \mbox{2023~--- 16}, \mbox{2024~--- 8}, \mbox{2025~--- 7}); \mbox{тестових~--- 2}; \mbox{на відповідність~--- 29}.}\par\vspace{0.2cm}

\typeTitle{Завдання з вибором однієї відповіді (А--Д)}

\begin{samepage}
\noindent\zadnum Графік функції $y = 3^x$ паралельно перенесли на 2 одиниці вниз уздовж осі $y$. Укажіть функцію, графік якої отримали в результаті цього перетворення. \nmtyear{2025}\par
\nopagebreak\nbvspace{0.3cm}

\nmtnobreak
\answerTableTall{$y = 3^{x-2}$}{$y = 3^x - 2$}{$y = 3^{x+2}$}{$y = 3^x + 2$}{$y = \displaystyle\frac{3^x}{2}$}
\par\penalty-20
\end{samepage}

\begin{samepage}
\noindent\zadnum Спростіть вираз $5^{1+x} \cdot 5^{1-x}$. \nmtyear{2025}\par
\nopagebreak\nbvspace{0.3cm}

\nmtnobreak
\answerTable{$25^{1-x^2}$}{$5$}{$25$}{$5^{1-x^2}$}{$10$}
\par\penalty-20
\end{samepage}

\typeTitle{Завдання на встановлення відповідності}

\begin{samepage}
\noindent\zadnum До кожного виразу (1–3) доберіть тотожно рівний йому вираз (А–Д), якщо $a \neq 8$. \nmtyear{2023}\par
\nopagebreak\nbvspace{0.3cm}

\nmtnobreak
\noindent
\matchingLayout{
\matchHead{Вираз}
    \matchItem{1}{$\displaystyle\frac{64 - a^2}{8 - a}$}
\matchItem{2}{$2^{3a}$}
\matchItem{3}{$\log_{\sqrt[3]{2}} 2^a$}
}{
\matchHead{Тотожно рівний вираз}
    \matchItem{А}{$8^a$}
\matchItem{Б}{$8 + a$}
\matchItem{В}{$8a$}
\matchItem{Г}{$8 - a$}
\matchItem{Д}{$\displaystyle\frac{a}{8}$}
}{
\matchingGrid
}
\par\penalty-20\vspace{0.25cm}
\end{samepage}

\begin{samepage}
\noindent\zadnum Установіть відповідність між виразом (1–3) і проміжком (А–Д), якому належить значення цього виразу. \nmtyear{2023}\par
\nopagebreak\nbvspace{0.3cm}

\nmtnobreak
\noindent
\matchingLayout{
\matchHead{Вираз}
    \matchItem{1}{$\cos \displaystyle\frac{\pi}{2}$}
\matchItem{2}{$|\pi - 5|$}
\matchItem{3}{$2^\pi$}
}{
\matchHead{Проміжок}
    \matchItem{А}{$(-\infty; 0]$}
\matchItem{Б}{$(0; 2)$}
\matchItem{В}{$[2; 4)$}
\matchItem{Г}{$[4; 8)$}
\matchItem{Д}{$[8; +\infty)$}
}{
\matchingGrid
}
\par\penalty-20\vspace{0.25cm}
\end{samepage}

\begin{samepage}
\noindent\zadnum До кожного початку речення (1–3) доберіть його закінчення (А–Д) так, щоб утворилося правильне твердження, якщо $m$ і $n$ --- натуральні числа, $n > 1$, $m > 1$. \nmtyear{2023}\par
\nopagebreak\nbvspace{0.3cm}

\nmtnobreak
\noindent
\matchingLayout{
\matchHead{Початок речення}
    \matchItem{1}{Якщо $n \sin m\pi = a$, то}
\matchItem{2}{Якщо $\displaystyle\frac{2^m}{2^n} = 2^a$, то}
\matchItem{3}{Якщо $\sqrt[n]{\sqrt[n]{2}} = \sqrt[4]{2}$, то}
}{
\matchHead{Закінчення речення}
    \matchItem{А}{$a = mn$.}
\matchItem{Б}{$a = 0$.}
\matchItem{В}{$a = m - n$.}
\matchItem{Г}{$a = n$.}
\matchItem{Д}{$a = \displaystyle\frac{m}{n}$.}
}{
\matchingGrid
}
\par\penalty-20\vspace{0.25cm}
\end{samepage}

\begin{samepage}
\noindent\zadnum Установіть відповідність між функцією (1–3) та її найменшим значенням (А–Д) на відрізку $[0; 4]$. \nmtyear{2023}\par
\nopagebreak\nbvspace{0.3cm}

\nmtnobreak
\noindent
\matchingLayout{
\matchHead{Функція}
    \matchItem{1}{$y = x^2 - 2x + 5$}
\matchItem{2}{$y = -0{,}5x + 5$}
\matchItem{3}{$y = 2^x$}
}{
\matchHead{Найменше значення на $[0; 4]$}
    \matchItem{А}{$1$}
\matchItem{Б}{$2$}
\matchItem{В}{$3$}
\matchItem{Г}{$4$}
\matchItem{Д}{$5$}
}{
\matchingGrid
}
\par\penalty-20\vspace{0.25cm}
\end{samepage}

\begin{samepage}
\noindent\zadnum До кожного початку речення (1–3) доберіть його закінчення (А–Д) так, щоб утворилося правильне твердження. \nmtyear{2023}\par
\nopagebreak\nbvspace{0.3cm}

\nmtnobreak
\noindent
\matchingLayout{
\matchHead{Початок речення}
    \matchItem{1}{Якщо $3^m \cdot 9^n = 3^a$, то}
\matchItem{2}{Якщо $(m + n)^2 - m^2 - n^2 = a$, то}
\matchItem{3}{Якщо $\log_3 27^{mn} = a$, то}
}{
\matchHead{Закінчення речення}
    \matchItem{А}{$a = 0$.}
\matchItem{Б}{$a = 27^{mn}$.}
\matchItem{В}{$a = 3mn$.}
\matchItem{Г}{$a = 2mn$.}
\matchItem{Д}{$a = m + 2n$.}
}{
\matchingGrid
}
\par\penalty-20\vspace{0.25cm}
\end{samepage}

\begin{samepage}
\noindent\zadnum \begin{minipage}[t]{0.55\textwidth}
У прямокутній системі координат на площині зображено коло $x^2+y^2=4$. Установіть відповідність між функцією (1–3) та кількістю (А – Д) спільних точок, які має графік цієї функції із заданим колом. \nmtyear{2023}
\end{minipage}
\hfill
\begin{minipage}[t]{0.4\textwidth}
    \nopagebreak\nbvspace{-0.3cm}
    \begin{flushright}
    \begin{nmtfit}\begin{tikzpicture}[scale=0.6]
        % Сітка та осі

        \draw[->, >=stealth, thick] (-3,0) -- (3,0) node[below] {$x$};
        \draw[->, >=stealth, thick] (0,-3) -- (0,3) node[left] {$y$};

        % Коло
        \draw[thick] (0,0) circle (2cm);

        % Підписи
        \node[below left] at (0,0) {$0$};
        \node[above right] at (1.4, 1.4) {$x^2+y^2=4$};
    \end{tikzpicture}\end{nmtfit}
    \end{flushright}
\end{minipage}

\nmtnobreak
\nopagebreak\nbvspace{0.3cm}

\nmtnobreak
\noindent
\matchingLayout{
\matchHead{Функція}
    \matchItem{1}{$y=x+4$}
\matchItem{2}{$y=x^2-2$}
\matchItem{3}{$y=4^x$}
}{
\matchHead{Кількість спільних точок}
    \matchItem{А}{жодної}
\matchItem{Б}{лише одна}
\matchItem{В}{лише дві}
\matchItem{Г}{лише три}
\matchItem{Д}{більше за три}
}{
\matchingGrid
}
\par\penalty-20\vspace{0.25cm}
\end{samepage}

\begin{samepage}
\noindent\zadnum Увідповідніть функцію (1–3) із кількістю (А–Д) спільних точок її графіка з прямою $y = -x$. \nmtyear{2023}\par
\nopagebreak\nbvspace{0.3cm}

\nmtnobreak
\noindent
\matchingLayout{
\matchHead{Функція}
    \matchItem{1}{$y = -\sqrt{x}$}
\matchItem{2}{$y = \left(\displaystyle\frac{1}{2}\right)^x$}
\matchItem{3}{$y = 2x + 2$}
}{
\matchHead{Кількість спільних точок}
    \matchItem{А}{жодної}
\matchItem{Б}{одна}
\matchItem{В}{дві}
\matchItem{Г}{три}
\matchItem{Д}{безліч}
}{
\matchingGrid
}
\par\penalty-20\vspace{0.25cm}
\end{samepage}

\begin{samepage}
\noindent\zadnum Установіть відповідність між функцією (1–3) та її найбільшим значенням (А–Д) на відрізку $[-1; 3]$. \nmtyear{2023}\par
\nopagebreak\nbvspace{0.3cm}

\nmtnobreak
\noindent
\matchingLayout{
\matchHead{Функція}
    \matchItem{1}{$y = x^2 - 2x$}
\matchItem{2}{$y = -x$}
\matchItem{3}{$y = 5^{-x}$}
}{
\matchHead{Найбільше значення на $[-1; 3]$}
    \matchItem{А}{$1$}
\matchItem{Б}{$2$}
\matchItem{В}{$3$}
\matchItem{Г}{$4$}
\matchItem{Д}{$5$}
}{
\matchingGrid
}
\par\penalty-20\vspace{0.25cm}
\end{samepage}

\begin{samepage}
\noindent\zadnum Увідповідніть функцію (1–3) із кількістю (А–Д) спільних точок її графіка з прямою $y = x + 3$. \nmtyear{2023}\par
\nopagebreak\nbvspace{0.3cm}

\nmtnobreak
\noindent
\matchingLayout{
\matchHead{Функція}
    \matchItem{1}{$y = x$}
\matchItem{2}{$y = 2^{-x}$}
\matchItem{3}{$y = \displaystyle\frac{1}{x}$}
}{
\matchHead{Кількість спільних точок}
    \matchItem{А}{жодної}
\matchItem{Б}{одна}
\matchItem{В}{дві}
\matchItem{Г}{три}
\matchItem{Д}{безліч}
}{
\matchingGrid
}
\par\penalty-20\vspace{0.25cm}
\end{samepage}

\begin{samepage}
\noindent\zadnum Установіть відповідність між функцією (1–3) та її властивістю (А–Д). \nmtyear{2023}\par
\nopagebreak\nbvspace{0.3cm}

\nmtnobreak
\noindent
\matchingLayout{
\matchHead{Функція}
    \matchItem{1}{$y = x^2 - 1$}
\matchItem{2}{$y = 3^x$}
\matchItem{3}{$y = x^3$}
}{
\matchHead{Властивість}
    \matchItem{А}{проходить через початок координат}
\matchItem{Б}{не містить точок з від'ємною ординатою}
\matchItem{В}{не має спільних точок з віссю $y$}
\matchItem{Г}{двічі перетинає графік прямої $y = 3$}
\matchItem{Д}{графік знаходиться лише в I чверті}
}{
\matchingGrid
}
\par\penalty-20\vspace{0.25cm}
\end{samepage}

\begin{samepage}
\noindent\zadnum До кожного початку речення (1–3) доберіть його закінчення (А–Д) так, щоб утворилося правильне твердження, якщо $n > 0$. \nmtyear{2023}\par
\nopagebreak\nbvspace{0.3cm}

\nmtnobreak
\noindent
\matchingLayout{
\matchHead{Початок речення}
    \matchItem{1}{Якщо $\displaystyle\frac{n}{a} = 3$, то}
\matchItem{2}{Якщо $1 + \log_3 n = \log_3 a$, то}
\matchItem{3}{Якщо $3^n \cdot 3 = 3^a$, то}
}{
\matchHead{Закінчення речення}
    \matchItem{А}{$a = 3n$.}
\matchItem{Б}{$a = n + 1$.}
\matchItem{В}{$a = n + 3$.}
\matchItem{Г}{$a = \displaystyle\frac{3}{n}$.}
\matchItem{Д}{$a = \displaystyle\frac{n}{3}$.}
}{
\matchingGrid
}
\par\penalty-20\vspace{0.25cm}
\end{samepage}

\begin{samepage}
\noindent\zadnum Установіть відповідність між твердженням (1–3) та функцією (А–Д), для якої це твердження є правильним. \nmtyear{2023}\par
\nopagebreak\nbvspace{0.3cm}

\nmtnobreak
\noindent
\matchingLayout{
\matchHead{Твердження}
    \matchItem{1}{областю визначення функції є проміжок $[-1; +\infty)$}
\matchItem{2}{графік функції проходить через точку $(0; 1)$}
\matchItem{3}{функція має точку локального екстремуму на проміжку $[1; 3]$}
}{
\matchHead{Функція}
    \matchItem{А}{$y = -\cos x$}
\matchItem{Б}{$y = 4^x$}
\matchItem{В}{$y = 2\sqrt{x + 1}$}
\matchItem{Г}{$y = x^2 - 4x + 3$}
\matchItem{Д}{$y = x$}
}{
\matchingGrid
}
\par\penalty-20\vspace{0.25cm}
\end{samepage}

\begin{samepage}
\noindent\zadnum Установіть відповідність між функцією (1–3) та її властивістю (А–Д). \nmtyear{2023}\par
\nopagebreak\nbvspace{0.3cm}

\nmtnobreak
\noindent
\matchingLayout{
\matchHead{Функція}
    \matchItem{1}{$y = (x + 2)^2$}
\matchItem{2}{$y = 2\sqrt{x}$}
\matchItem{3}{$y = 2^x$}
}{
\matchHead{Властивість}
    \matchItem{А}{є зростаючою на $(-\infty; +\infty)$}
\matchItem{Б}{графік має одну спільну точку із графіком $y = x - 2$}
\matchItem{В}{графік має дві спільні точки із графіком $y = x - 2$}
\matchItem{Г}{є спадною на проміжку $(-\infty; -2]$}
\matchItem{Д}{є спадною на проміжку $(-\infty; 2]$}
}{
\matchingGrid
}
\par\penalty-20\vspace{0.25cm}
\end{samepage}

\begin{samepage}
\noindent\zadnum Установіть відповідність між виразом (1–3) і множиною (А–Д), якій належить значення цього виразу, якщо $a = 5$. \nmtyear{2023}\par
\nopagebreak\nbvspace{0.3cm}

\nmtnobreak
\noindent
\matchingLayout{
\matchHead{Вираз}
    \matchItem{1}{$\log_{0{,}2} a$}
\matchItem{2}{$2^{-a}$}
\matchItem{3}{$\displaystyle\frac{10}{a}$}
}{
\matchHead{Множина}
    \matchItem{А}{ірраціональних додатних чисел}
\matchItem{Б}{натуральних чисел}
\matchItem{В}{цілих від'ємних чисел}
\matchItem{Г}{раціональних нецілих чисел}
\matchItem{Д}{ірраціональних від'ємних чисел}
}{
\matchingGrid
}
\par\penalty-20\vspace{0.25cm}
\end{samepage}

\begin{samepage}
\noindent\zadnum До кожного початку речення (1–3) доберіть його закінчення (А–Д) так, щоб утворилося правильне твердження. \nmtyear{2023}\par
\nopagebreak\nbvspace{0.3cm}

\nmtnobreak
\noindent
\matchingLayout{
\matchHead{Початок речення}
    \matchItem{1}{Функція $y = \sqrt{x + 1}$}
\matchItem{2}{Функція $y = 4 - x^2$}
\matchItem{3}{Функція $y = 3^{-x}$}
}{
\matchHead{Закінчення речення}
    \matchItem{А}{має точку локального максимуму.}
\matchItem{Б}{має точку локального мінімуму.}
\matchItem{В}{є непарною.}
\matchItem{Г}{зростає на всій області визначення.}
\matchItem{Д}{набуває лише додатних значень.}
}{
\matchingGrid
}
\par\penalty-20\vspace{0.25cm}
\end{samepage}

\begin{samepage}
\noindent\zadnum Установіть відповідність між функцією (1–3) та властивістю її графіка (А–Д). \nmtyear{2023}\par
\nopagebreak\nbvspace{0.3cm}

\nmtnobreak
\noindent
\matchingLayout{
\matchHead{Функція}
    \matchItem{1}{$y = \cos x$}
\matchItem{2}{$y = -\displaystyle\frac{2}{x}$}
\matchItem{3}{$y = 2^x$}
}{
\matchHead{Властивість графіка}
    \matchItem{А}{проходить через точку $(1; 0)$}
\matchItem{Б}{не перетинає вісь $y$}
\matchItem{В}{симетричний відносно осі $y$}
\matchItem{Г}{розміщений лише в I та II чвертях}
\matchItem{Д}{розміщений лише в I та IV чвертях}
}{
\matchingGrid
}
\par\penalty-20\vspace{0.25cm}
\end{samepage}

\begin{samepage}
\noindent\zadnum Установіть відповідність між твердженням (1–3) та функцією (А–Д), для якої це твердження є правильним. \nmtyear{2024}\par
\nopagebreak\nbvspace{0.3cm}

\nmtnobreak
\noindent
\matchingLayout{
\matchHead{Твердження}
    \matchItem{1}{функція має 2 нулі}
\matchItem{2}{на відрізку $[-1; 3]$ функція набуває від'ємних значень}
\matchItem{3}{найменше значення функції на відрізку $[-1; 3]$ дорівнює $0{,}5$}
}{
\matchHead{Функція}
    \matchItem{А}{$y = x^2 - 4$}
\matchItem{Б}{$y = \displaystyle\frac{1}{x - 4}$}
\matchItem{В}{$y = 2^x$}
\matchItem{Г}{$y = 0{,}5^x$}
\matchItem{Д}{$y = \sqrt{x + 1}$}
}{
\matchingGrid
}
\par\penalty-20\vspace{0.25cm}
\end{samepage}

\begin{samepage}
\noindent\zadnum Установіть відповідність між функцією (1–3) та її властивістю (А–Д). \nmtyear{2024}\par
\nopagebreak\nbvspace{0.3cm}

\nmtnobreak
\noindent
\matchingLayout{
\matchHead{Функція}
    \matchItem{1}{$y = 4 - x^2$}
\matchItem{2}{$y = -x^3$}
\matchItem{3}{$y = 4^x$}
}{
\matchHead{Властивість}
    \matchItem{А}{є зростаючою на всій області визначення}
\matchItem{Б}{набуває від'ємного значення при $x = -2$}
\matchItem{В}{є парною}
\matchItem{Г}{має три спільні точки з графіком функції $y = -x$}
\matchItem{Д}{графік функції розташований лише в I та IV координатній чверті}
}{
\matchingGrid
}
\par\penalty-20\vspace{0.25cm}
\end{samepage}

\begin{samepage}
\noindent\zadnum Установіть відповідність між виразом (1–3) та проміжком (А–Д), якому належить значення цього виразу. \nmtyear{2024}\par
\nopagebreak\nbvspace{0.3cm}

\nmtnobreak
\noindent
\matchingLayout{
\matchHead{Вираз}
    \matchItem{1}{$\mathrm{tg}\, \displaystyle\frac{\pi}{3}$}
\matchItem{2}{$1 - \pi$}
\matchItem{3}{$\left(\displaystyle\frac{1}{2}\right)^{\log_2 \pi}$}
}{
\matchHead{Проміжок}
    \matchItem{А}{$[-5; -2)$}
\matchItem{Б}{$[-2; 0)$}
\matchItem{В}{$[0; 1)$}
\matchItem{Г}{$[1; 2)$}
\matchItem{Д}{$[2; 5)$}
}{
\matchingGrid
}
\par\penalty-20\vspace{0.25cm}
\end{samepage}

\begin{samepage}
\noindent\zadnum Установіть відповідність між функцією (1–3) та властивістю її графіка (А–Д). \nmtyear{2024}\par
\nopagebreak\nbvspace{0.3cm}

\nmtnobreak
\noindent
\matchingLayout{
\matchHead{Функція}
    \matchItem{1}{$y = \sqrt{x + 4}$}
\matchItem{2}{$y = \displaystyle\frac{4}{x}$}
\matchItem{3}{$y = 4^x$}
}{
\matchHead{Властивість графіка функції}
    \matchItem{А}{розташований лише в першій чверті}
\matchItem{Б}{має спільну точку з віссю $x$}
\matchItem{В}{проходить через точку $(0; 1)$}
\matchItem{Г}{не перетинає вісь $y$}
\matchItem{Д}{симетричний відносно осі $y$}
}{
\matchingGrid
}
\par\penalty-20\vspace{0.25cm}
\end{samepage}

\begin{samepage}
\noindent\zadnum Доберіть до початку речення (1–3) його закінчення (А–Д) так, щоб утворилося правильне твердження. \nmtyear{2024}\par
\nopagebreak\nbvspace{0.3cm}

\nmtnobreak
\noindent
\matchingLayout{
\matchHead{Початок речення}
    \matchItem{1}{Графік функції $y = x$ не має жодної спільної точки з}
\matchItem{2}{Графік функції $y = 3^x$ має лише одну спільну точку з}
\matchItem{3}{Графік рівняння $(x + 3)^2 + y^2 = 4$ має дві спільні точки з}
}{
\matchHead{Закінчення речення}
    \matchItem{А}{віссю $x$.}
\matchItem{Б}{віссю $y$.}
\matchItem{В}{прямою $y = x - 4$.}
\matchItem{Г}{прямою $y = -4$.}
\matchItem{Д}{прямою $y = -2$.}
}{
\matchingGrid
}
\par\penalty-20\vspace{0.25cm}
\end{samepage}

\begin{samepage}
\noindent\zadnum До кожного початку речення (1–3) доберіть його закінчення (А–Д) так, щоб утворилося правильне твердження. \nmtyear{2024}\par
\nopagebreak\nbvspace{0.3cm}

\nmtnobreak
\noindent
\matchingLayout{
\matchHead{Початок речення}
    \matchItem{1}{Графік функції $y = 2x$}
\matchItem{2}{Графік функції $y = \log_2 x$}
\matchItem{3}{Графік функції $y = 2^x$}
}{
\matchHead{Закінчення речення}
    \matchItem{А}{симетричний відносно осі абсцис.}
\matchItem{Б}{симетричний відносно осі ординат.}
\matchItem{В}{симетричний відносно початку координат.}
\matchItem{Г}{не перетинає вісь абсцис.}
\matchItem{Д}{не перетинає вісь ординат.}
}{
\matchingGrid
}
\par\penalty-20\vspace{0.25cm}
\end{samepage}

\begin{samepage}
% рік визначено за сусідніми завданнями (у старому файлі тег року відсутній)
\noindent\zadnum На кожному з рисунків \mbox{(1--3)} зображено графік функції $y=f(x)$, визначеної на проміжку $[-3; 3]$. Узгодьте рисунок \mbox{(1--3)} з твердженням \mbox{(А--Д)} щодо функції $y=f(x)$, графік якої зображено на цьому рисунку. \nmtyear{2024}

\nmtnobreak
\nopagebreak\nbvspace{0.3cm}

\nmtnobreak
% Блок рисунків
\noindent
\begin{minipage}[t]{0.32\textwidth}\nbvspace{0pt}
    \centering
    Рис. 1 \\
    \begin{nmtfit}\begin{tikzpicture}[scale=0.45]
        \draw[step=1cm,gray!50,very thin] (-3.5,-3.5) grid (3.5,3.5);
        \draw[->, >=stealth] (-3.5,0) -- (3.5,0) node[below] {$x$};
        \draw[->, >=stealth] (0,-3.5) -- (0,3.5) node[left] {$y$};
        \node[below left] at (0,0) {\small $0$};
        \node[left] at (0,1) {\small $1$};
        \node[below] at (1,0) {\small $1$};
        \node[below] at (3,0) {\small $3$};
        \node[below] at (-3,0) {\small $-3$};

        % Дуга (схожа на півокружність або параболу вниз)
        \draw[thick] plot [smooth, tension=1.5] coordinates {(-3, -1) (0, -4) (3, -1)};
        \node[right] at (2, -2.5) {\small $y=f(x)$};
        \fill (-3,-1) circle (3pt);
        \fill (3,-1) circle (3pt);
    \end{tikzpicture}\end{nmtfit}
\end{minipage}
\hfill
\begin{minipage}[t]{0.32\textwidth}\nbvspace{0pt}
    \centering
    Рис. 2 \\
    \begin{nmtfit}\begin{tikzpicture}[scale=0.45]
        \draw[step=1cm,gray!50,very thin] (-3.5,-3.5) grid (3.5,3.5);
        \draw[->, >=stealth] (-3.5,0) -- (3.5,0) node[below] {$x$};
        \draw[->, >=stealth] (0,-3.5) -- (0,3.5) node[left] {$y$};
        \node[below left] at (0,0) {\small $0$};
        \node[left] at (0,1) {\small $1$};
        \node[below] at (1,0) {\small $1$};
        \node[below] at (3,0) {\small $3$};
        \node[below] at (-3,0) {\small $-3$};

        % "Горб"
        \draw[thick] plot [smooth, tension=0.6] coordinates {(-3, -2) (-1, 3) (1, 0) (3, -1.5)};
        \node[right] at (1, 2) {\small $y=f(x)$};
        \fill (-3,-2) circle (3pt);
        \fill (3,-1.5) circle (3pt);
    \end{tikzpicture}\end{nmtfit}
\end{minipage}
\hfill
\begin{minipage}[t]{0.32\textwidth}\nbvspace{0pt}
    \centering
    Рис. 3 \\
    \begin{nmtfit}\begin{tikzpicture}[scale=0.45]
        \draw[step=1cm,gray!50,very thin] (-3.5,-3.5) grid (3.5,3.5);
        \draw[->, >=stealth] (-3.5,0) -- (3.5,0) node[below] {$x$};
        \draw[->, >=stealth] (0,-3.5) -- (0,3.5) node[left] {$y$};
        \node[below left] at (0,0) {\small $0$};
        \node[left] at (0,1) {\small $1$};
        \node[below] at (1,0) {\small $1$};
        \node[below] at (3,0) {\small $3$};
        \node[below] at (-3,0) {\small $-3$};

        % Спадна функція
        \draw[thick] plot [smooth, tension=0.6] coordinates {(-2.5, 3.5) (-1, 2) (0, 0) (1, -1) (3, -3)};
        \node[left] at (-2, 2) {\small $y=f(x)$};
        \fill (3,-3) circle (3pt);
    \end{tikzpicture}\end{nmtfit}
\end{minipage}

\nmtnobreak
\nopagebreak\nbvspace{0.3cm}

\nmtnobreak
\noindent
\matchingLayout{
\matchHead{Рисунок}
\matchItem{1}{Рис. 1}
\matchItem{2}{Рис. 2}
\matchItem{3}{Рис. 3}
}{
\matchHead{Твердження}
    \matchItem{А}{графік функції $f$ є фрагментом кола $x^2 + (y-1)^2 = 9$}
\matchItem{Б}{графік функції $f$ є фрагментом кола $x^2 + (y+1)^2 = 9$}
\matchItem{В}{функція $f$ зростає на області визначення}
\matchItem{Г}{графік функції $y=f(x)+2$ розташований лише в I і II чвертях}
\matchItem{Д}{графік функції $f$ має лише одну спільну точку з графіком функції $y=2^x$}
}{
\matchingGrid
}
\par\penalty-20\vspace{0.25cm}
\end{samepage}

\begin{samepage}
% рік визначено за сусідніми завданнями (у старому файлі тег року відсутній)
\noindent\zadnum Установіть відповідність між твердженням \mbox{(1--3)} та прямою, зображеною на рисунку \mbox{(А--Д)}, для якої це твердження є правильним. \nmtyear{2024}

\nmtnobreak
\nopagebreak\nbvspace{0.3cm}

\nmtnobreak
\noindent
\begin{minipage}[t]{0.60\textwidth}
\matchHead{Твердження}
    \matchItem{1}{не має спільних точок з функцією $y=\left(\dfrac{2}{3}\right)^x$}
\matchItem{2}{є графіком функції $y=x-2$}
\matchItem{3}{кутовий коефіцієнт прямої дорівнює $0$}
\end{minipage}%
\hfill
\begin{minipage}[t]{0.30\textwidth}
    \nopagebreak\nbvspace{0.5cm}
    \matchingGrid
\end{minipage}

\nmtnobreak
\nopagebreak\nbvspace{0.5cm}
\matchHead{Пряма}

\nmtnobreak
% Таблиця з графіками
\begin{center}
\begingroup
\setlength{\tabcolsep}{3pt}
\begin{tabular}{|*{5}{c|}}
\hline
\textbf{А} & \textbf{Б} & \textbf{В} & \textbf{Г} & \textbf{Д} \\
\hline
\begin{nmtfit}\begin{tikzpicture}[scale=0.35]
    \draw[->, >=stealth] (-1.5,0) -- (3.5,0) node[below] {$x$};
    \draw[->, >=stealth] (0,-3) -- (0,2) node[left] {$y$};
    \node[below left] at (0,0) {\tiny $0$};
    \node[below] at (2,0) {\tiny $2$};
    \node[left] at (0,-2) {\tiny $-2$};
    \draw[thick] (-0.5,-2.5) -- (2.5,0.5); % y = x - 2
\end{tikzpicture}\end{nmtfit} &
\begin{nmtfit}\begin{tikzpicture}[scale=0.35]
    \draw[->, >=stealth] (-2.5,0) -- (2.5,0) node[below] {$x$};
    \draw[->, >=stealth] (0,-2) -- (0,3) node[left] {$y$};
    \node[below left] at (0,0) {\tiny $0$};
    \node[above right] at (0,1) {\tiny $1$};
    \draw (0.1,1) -- (0,1) -- (0.2,1.2) -- (0.2,1); % прямий кут
    \draw[thick] (-2.5,1) -- (2.5,1); % y = 1
\end{tikzpicture}\end{nmtfit} &
\begin{nmtfit}\begin{tikzpicture}[scale=0.35]
    \draw[->, >=stealth] (-2.5,0) -- (2.5,0) node[below] {$x$};
    \draw[->, >=stealth] (0,-2) -- (0,3) node[left] {$y$};
    \node[below right] at (0,0) {\tiny $0$};
    \node[below left] at (-1,0) {\tiny $-1$};
    \draw (-1,0) -- (-0.8,0) -- (-0.8,0.2) -- (-1,0.2); % прямий кут
    \draw[thick] (-1,-2) -- (-1,3); % x = -1
\end{tikzpicture}\end{nmtfit} &
\begin{nmtfit}\begin{tikzpicture}[scale=0.35]
    \draw[->, >=stealth] (-3.5,0) -- (2.5,0) node[below] {$x$};
    \draw[->, >=stealth] (0,-3) -- (0,2) node[left] {$y$};
    \node[below right] at (0,0) {\tiny $0$};
    \node[below] at (-2,0) {\tiny $-2$};
    \node[left] at (0,-2) {\tiny $-2$};
    \draw[thick] (-2.5,0.5) -- (0.5,-2.5); % y = -x - 2
\end{tikzpicture}\end{nmtfit} &
\begin{nmtfit}\begin{tikzpicture}[scale=0.35]
    \draw[->, >=stealth] (-1.5,0) -- (4,0) node[below] {$x$};
    \draw[->, >=stealth] (0,-3) -- (0,3) node[left] {$y$};
    \node[below left] at (0,0) {\tiny $0$};
    \node[below] at (3,0) {\tiny $3$};
    \node[left] at (0,-2) {\tiny $-2$};
    \draw[thick] (-1,-2.66) -- (4,0.66); % y = 2/3 x - 2
\end{tikzpicture}\end{nmtfit} \\
\hline
\end{tabular}
\endgroup
\end{center}
\par\penalty-20\vspace{0.25cm}
\end{samepage}

\begin{samepage}
\noindent\zadnum Узгодьте функцію (1–3) із кількістю (А–Д) спільних точок її графіка з прямою $y = -x$. \nmtyear{2025}\par
\nopagebreak\nbvspace{0.3cm}

\nmtnobreak
\noindent
\matchingLayout{
\matchHead{Функція}
    \matchItem{1}{$y = -x^3$}
\matchItem{2}{$y = 2 - x$}
\matchItem{3}{$y = 2^x$}
}{
\matchHead{Кількість спільних точок}
    \matchItem{А}{жодної}
\matchItem{Б}{одна}
\matchItem{В}{дві}
\matchItem{Г}{три}
\matchItem{Д}{безліч}
}{
\matchingGrid
}
\par\penalty-20\vspace{0.25cm}
\end{samepage}

\begin{samepage}
\noindent\zadnum Узгодьте вираз (1–3) із значенням $m$ (А–Д), за якого значення цього виразу дорівнює 1. \nmtyear{2025}\par
\nopagebreak\nbvspace{0.3cm}

\nmtnobreak
\noindent
\matchingLayout{
\matchHead{Вираз}
    \matchItem{1}{$\displaystyle\frac{m}{4}$}
\matchItem{2}{$4^6 : 4^{-m}$}
\matchItem{3}{$\log_{16} 2 + \log_{16} m$}
}{
\matchHead{Значення $m$}
    \matchItem{А}{$\displaystyle\frac{1}{4}$}
\matchItem{Б}{$6$}
\matchItem{В}{$4$}
\matchItem{Г}{$-6$}
\matchItem{Д}{$8$}
}{
\matchingGrid
}
\par\penalty-20\vspace{0.25cm}
\end{samepage}

\begin{samepage}
\noindent\zadnum Доберіть до кожного початку речення (1–3) його закінчення (А–Д) так, щоб утворилося правильне твердження, якщо $n \neq 0$. \nmtyear{2025}\par
\nopagebreak\nbvspace{0.3cm}

\nmtnobreak
\noindent
\matchingLayout{
\matchHead{Початок речення}
    \matchItem{1}{Якщо $(m + 1)^2 - 1 = n$, то}
\matchItem{2}{Якщо $\sqrt{2^m} = 2^n$, то}
\matchItem{3}{Якщо $\log_{2^n} 4^m = 1$, то}
}{
\matchHead{Закінчення речення}
    \matchItem{А}{$n = 2m$.}
\matchItem{Б}{$n = m^2$.}
\matchItem{В}{$n = m^2 + 2m$.}
\matchItem{Г}{$n = \sqrt{m}$.}
\matchItem{Д}{$n = \displaystyle\frac{m}{2}$.}
}{
\matchingGrid
}
\par\penalty-20\vspace{0.25cm}
\end{samepage}

\begin{samepage}
% рік визначено за сусідніми завданнями (у старому файлі тег року відсутній)
\noindent\zadnum \begin{minipage}[t]{0.55\textwidth}
На рисунку зображено графік функції $y=f(x)$, визначеної на проміжку $[-4; 6]$. Узгодьте проміжок \mbox{(1--3)} із функцією \mbox{(А--Д)}, що описує графік функції $y=f(x)$ на цьому проміжку. \nmtyear{2025}
\end{minipage}
\hfill
\begin{minipage}[t]{0.4\textwidth}
    \nopagebreak\nbvspace{-0.3cm}
    \begin{flushright}
    \begin{nmtfit}\begin{tikzpicture}[scale=0.5]
        \draw[step=1cm,gray!50,very thin] (-4.5,-0.5) grid (6.5,8.5);
        \draw[->, >=stealth, thick] (-4.5,0) -- (6.5,0) node[below] {$x$};
        \draw[->, >=stealth, thick] (0,-0.5) -- (0,8.5) node[left] {$y$};

        \node[below left] at (0,0) {$0$};
        \node[below] at (1,0) {$1$};
        \node[left] at (0,1) {$-1$}; % На картинці тут -1 позначено як засічка, але насправді це 1, просто риска
        \node[left] at (0,1) {$1$};
        \node[below] at (6,0) {$6$};
        \node[below] at (-4,0) {$-4$};

        % Графік
        % [-4; 0] -> y = sqrt(-x). (-4, 2) -> (0,0)
        \draw[thick] plot [smooth, domain=-4:0] (\x, {sqrt(-\x)});

        % [0; 2] -> y = x^3. (0,0) -> (2, 8)
        \draw[thick] plot [smooth, domain=0:2] (\x, {\x*\x*\x});

        % [2; 6] -> y = 12 - 2x. (2, 8) -> (6, 0)
        \draw[thick] (2, 8) -- (6, 0);

        \fill (-4, 2) circle (3pt);
        \fill (6, 0) circle (3pt);
        \node[right, fill=white, inner sep=1.5pt] at (4, 4) {$y=f(x)$};
    \end{tikzpicture}\end{nmtfit}
    \end{flushright}
\end{minipage}

\nmtnobreak
\nopagebreak\nbvspace{0.3cm}

\nmtnobreak
\noindent
\matchingLayout{
\matchHead{Проміжок}
    \matchItem{1}{$[-4; 0]$}
\matchItem{2}{$(0; 2]$}
\matchItem{3}{$(2; 6]$}
}{
\matchHead{Функція}
    \matchItem{А}{$y=x^3$}
\matchItem{Б}{$y=12-2x$}
\matchItem{В}{$y=-\sqrt{x}$}
\matchItem{Г}{$y=2^x$}
\matchItem{Д}{$y=\sqrt{-x}$}
}{
\matchingGrid
}
\par\penalty-20\vspace{0.25cm}
\end{samepage}

\begin{samepage}
% рік визначено за сусідніми завданнями (у старому файлі тег року відсутній)
\noindent\zadnum Установіть відповідність між твердженням \mbox{(1--3)} та прямою, зображеною на рисунку \mbox{(А--Д)}, для якої це твердження є правильним. \nmtyear{2025}

\nmtnobreak
\nopagebreak\nbvspace{0.3cm}

\nmtnobreak
\noindent
\begin{minipage}[t]{0.60\textwidth}
\matchHead{Твердження}
    \matchItem{1}{не має спільних точок з функцією $y=\log_2(x-1)$}
\matchItem{2}{є графіком функції $y=\dfrac{2x}{3}-2$}
\matchItem{3}{кутовий коефіцієнт прямої є від'ємним числом}
\end{minipage}%
\hfill
\begin{minipage}[t]{0.30\textwidth}
    \nopagebreak\nbvspace{0.5cm}
    \matchingGrid
\end{minipage}

\nmtnobreak
\nopagebreak\nbvspace{0.5cm}
\matchHead{Пряма}

\nmtnobreak
% Таблиця з графіками
\begin{center}
\begingroup
\setlength{\tabcolsep}{3pt}
\begin{tabular}{|*{5}{c|}}
\hline
\textbf{А} & \textbf{Б} & \textbf{В} & \textbf{Г} & \textbf{Д} \\
\hline
\begin{nmtfit}\begin{tikzpicture}[scale=0.35]
    \draw[->, >=stealth] (-1.5,0) -- (3.5,0) node[below] {$x$};
    \draw[->, >=stealth] (0,-3) -- (0,2) node[left] {$y$};
    \node[below left] at (0,0) {\tiny $0$};
    \node[below] at (2,0) {\tiny $2$};
    \node[left] at (0,-2) {\tiny $-2$};
    \draw[thick] (-0.5,-2.5) -- (2.5,0.5); % y = x - 2
\end{tikzpicture}\end{nmtfit} &
\begin{nmtfit}\begin{tikzpicture}[scale=0.35]
    \draw[->, >=stealth] (-2.5,0) -- (2.5,0) node[below] {$x$};
    \draw[->, >=stealth] (0,-2) -- (0,3) node[left] {$y$};
    \node[below left] at (0,0) {\tiny $0$};
    \node[above right] at (0,1) {\tiny $1$};
    \draw (0.1,1) -- (0,1) -- (0.2,1.2) -- (0.2,1); % прямий кут
    \draw[thick] (-2.5,1) -- (2.5,1); % y = 1
\end{tikzpicture}\end{nmtfit} &
\begin{nmtfit}\begin{tikzpicture}[scale=0.35]
    \draw[->, >=stealth] (-2.5,0) -- (2.5,0) node[below] {$x$};
    \draw[->, >=stealth] (0,-2) -- (0,3) node[left] {$y$};
    \node[below right] at (0,0) {\tiny $0$};
    \node[below left] at (-1,0) {\tiny $-1$};
    \draw (-1,0) -- (-0.8,0) -- (-0.8,0.2) -- (-1,0.2); % прямий кут
    \draw[thick] (-1,-2) -- (-1,3); % x = -1
\end{tikzpicture}\end{nmtfit} &
\begin{nmtfit}\begin{tikzpicture}[scale=0.35]
    \draw[->, >=stealth] (-3.5,0) -- (2.5,0) node[below] {$x$};
    \draw[->, >=stealth] (0,-3) -- (0,2) node[left] {$y$};
    \node[below right] at (0,0) {\tiny $0$};
    \node[below] at (-2,0) {\tiny $-2$};
    \node[left] at (0,-2) {\tiny $-2$};
    \draw[thick] (-2.5,0.5) -- (0.5,-2.5); % y = -x - 2
\end{tikzpicture}\end{nmtfit} &
\begin{nmtfit}\begin{tikzpicture}[scale=0.35]
    \draw[->, >=stealth] (-1.5,0) -- (4,0) node[below] {$x$};
    \draw[->, >=stealth] (0,-3) -- (0,3) node[left] {$y$};
    \node[below left] at (0,0) {\tiny $0$};
    \node[below] at (3,0) {\tiny $3$};
    \node[left] at (0,-2) {\tiny $-2$};
    \draw[thick] (-1,-2.66) -- (4,0.66); % y = 2/3 x - 2
\end{tikzpicture}\end{nmtfit} \\
\hline
\end{tabular}
\endgroup
\end{center}
\par\penalty-20\vspace{0.25cm}
\end{samepage}

\endgroup
\end{document}
